\documentclass[12pt,a4paper]{article}

% ============================================================
% AI+Security 论文模板
% 作者: To-be-w1th0ut
% 描述: 用于撰写 AI 与安全交叉领域学术论文的 LaTeX 模板
% ============================================================

% --- 基础宏包 ---
\usepackage[utf8]{inputenc}
\usepackage[T1]{fontenc}
\usepackage{amsmath, amssymb, amsthm}
\usepackage{graphicx}
\usepackage{booktabs}
\usepackage{hyperref}
\usepackage{cleveref}
\usepackage{algorithm}
\usepackage{algpseudocode}
\usepackage{cite}

% --- 页面设置 ---
\usepackage[margin=2.5cm]{geometry}
\usepackage{setspace}
\onehalfspacing

% --- 颜色支持 ---
\usepackage{xcolor}

% --- 定理环境 ---
\newtheorem{theorem}{Theorem}[section]
\newtheorem{lemma}[theorem]{Lemma}
\newtheorem{definition}[theorem]{Definition}
\newtheorem{proposition}[theorem]{Proposition}

% --- 文档信息 ---
\title{\textbf{Your Paper Title Here}\\
\large A Study on AI and Security}
\author{Your Name\\
\textit{Your Institution}\\
\texttt{your.email@example.com}}
\date{\today}

\begin{document}

\maketitle

\begin{abstract}
在此输入你的摘要。摘要应简洁地概括研究的目的、方法、主要结果和结论。
通常 150-250 字为宜。

\medskip
\noindent\textbf{Keywords:} Artificial Intelligence, Cybersecurity, Machine Learning, Deep Learning, Threat Detection
\end{abstract}

\section{Introduction}
\label{sec:introduction}

在此介绍研究背景、动机和主要贡献。

\subsection{Background}
介绍 AI 与安全交叉领域的研究背景。

\subsection{Motivation}
阐述本研究的动机和现实意义。

\subsection{Contributions}
列出本文的主要贡献：
\begin{itemize}
    \item 贡献一
    \item 贡献二
    \item 贡献三
\end{itemize}

\section{Related Work}
\label{sec:related_work}

综述相关研究工作，包括：
\begin{itemize}
    \item AI 在安全领域的应用现状
    \item 传统安全方法的局限性
    \item 现有研究的不足
\end{itemize}

\section{Methodology}
\label{sec:methodology}

\subsection{Problem Formulation}
形式化描述研究问题。

\subsection{Proposed Approach}
详细介绍你提出的方法。

\begin{algorithm}[t]
\caption{Your Algorithm Name}
\label{alg:main}
\begin{algorithmic}[1]
\State \textbf{Input:} Training data $\mathcal{D}$, model parameters $\theta$
\State \textbf{Output:} Trained model $f_\theta$
\While{not converged}
    \State Sample batch $\mathcal{B} \sim \mathcal{D}$
    \State Compute gradient $g = \nabla_\theta \mathcal{L}(\mathcal{B}; \theta)$
    \State Update $\theta \leftarrow \theta - \eta g$
\EndWhile
\State \Return $f_\theta$
\end{algorithmic}
\end{algorithm}

\section{Experiments}
\label{sec:experiments}

\subsection{Experimental Setup}
描述实验环境、数据集和评估指标。

\subsection{Results}
展示实验结果。

\begin{table}[t]
\centering
\caption{Experimental Results Comparison}
\label{tab:results}
\begin{tabular}{lccc}
\toprule
\textbf{Method} & \textbf{Accuracy} & \textbf{Precision} & \textbf{Recall} \\
\midrule
Baseline & 0.85 & 0.83 & 0.82 \\
Method A & 0.89 & 0.88 & 0.87 \\
\textbf{Ours} & \textbf{0.94} & \textbf{0.93} & \textbf{0.92} \\
\bottomrule
\end{tabular}
\end{table}

\subsection{Analysis}
分析实验结果，解释原因。

\section{Discussion}
\label{sec:discussion}

讨论方法的优势、局限性和未来改进方向。

\section{Conclusion}
\label{sec:conclusion}

总结全文，重申主要贡献和意义。

% --- 参考文献 ---
\bibliographystyle{IEEEtran}
\bibliography{references}

\end{document}
